\section{CONFIGURING WORKFLOW}\label{sec:workflow}

\subsection{Basic Pattern}

Our modular structure supports practical configuration and improves reproducibility across heterogeneous XR deployments.
{\bf Figure~\ref{fig:authoring}} illustrates the configuring workflow for a basic two-hand and upper body setup.
Users first instantiate a Trigger Manager, register a Tracking Manager as a trigger, and configure the Tracking Manager’s device handlers. 
For upper-body motion, they insert an upper-body regression node and connect the Tracking Manager output as its input through sender--receiver interface matching.
They then bind retargeter nodes to both hand-tracking and upper-body outputs, and finally connect the retargeting stage to the target Motion Avatar format.
Users can validate this process by iterating between Editor and Play mode. 
This configured pipeline can be saved as the reusable preset.

Through this uniform workflow, users do not need to manually resolve SDK interface differences when binding motion to a custom avatar. This consistent pattern supports the tracking devices listed in \textbf{Table~\ref{tab:sdk-support-list}}.

\subsection{Extending the Authoring Pattern}
The same authoring pattern extends to richer pipelines by inserting additional computational nodes between tracking and retargeting stages while preserving the same runtime contracts.
{\bf Figure~\ref{fig:teaser}} shows representative pipeline configurations across heterogeneous XR devices for full-body avatar motion.
With Meta SDK full-body tracking, the graph can be configured as a relatively direct path from tracking to retargeting.
For resource-constrained devices, such as XREAL, selected computational nodes can be offloaded through network wrappers.
With Vive physical tracker inputs, the graph can be composed with IK nodes to exploit richer tracking signals.


